% \subsection{Sequences}
\tsDef{(Sequence)}{
    A sequence $(a_n)_{n \in \mathbb{N}_0}$ (also written $(a_n)_{n \ge 0}$ or $(a_n)_{n=0}^{\infty}$)
    is an infinite list of numbers, where each entry is called a \emph{term} of the sequence.
    A sequence can also be viewed as a function
    \(
        f : \mathbb{N}_0 \to \mathbb{R}
        \quad (\text{or } f : \mathbb{N} \to \mathbb{R}),
    \)
    whose values are
    \(
    f(n) = a_n.
    \)
}
\\
\tsDef{(Ways to Describe a Sequence)}{
    One way to describe a sequence is by giving an explicit formula for each term:
    \(
        a_n = f(n).
    \)
    Another possibility is to define the sequence recursively, constructing each new term
    from previous ones:
    \(
        a_n = g(a_{n-1}, a_{n-2}, \dots).
    \)
    The functions $f$ and $g$ are called the \emph{rule of formation} of the sequence.
}
\\
\tsDef{(Convergence/Divergence)}{
    A sequence $(a_n)_{n \in \mathbb{N}_0} \subset \mathbb{R}$ is called \emph{convergent}
    if there exists a number $L \in \mathbb{R}$ such that
    \[
        \forall \varepsilon > 0 \; \exists N > 0 \; \text{such that} \;
        \forall n > N : |a_n - L| < \varepsilon .
    \]
    The number $L$ is called the \emph{limit} of the sequence, and we write
    \[
        \lim_{n \to \infty} a_n = L.
    \]
    If no such $L$ exists, the sequence is called \emph{divergent}.
}
\\
\tsThe{(Uniqueness of the Limit)}{
    A convergent sequence has exactly one unique limit.
}
\\
\tsDef{(Divergence to infinity)}{
    Let $(a_n)_{n \in \mathbb{N}_0}$ be a sequence.

    \begin{itemize}
        \item If
              \(
                  \forall M > 0 \; \exists N > 0 \; \text{such that} \;
                  \forall n > N : a_n > M,
              \)
              then the sequence diverges to infinity and we write
              \(
                  \lim_{n \to \infty} a_n = \infty.
              \)
        \item If
              \(
                  \forall M > 0 \; \exists N > 0 \; \text{such that} \;
                  \forall n > N : a_n < -M,
              \)
              then the sequence diverges to minus infinity and we write
              \(
                  \lim_{n \to \infty} a_n = -\infty.
              \)
    \end{itemize}
}
\tsDef{(Subsequence)}{
    Let $(a_n)_{n \in \mathbb{N}_0}$ be a sequence in $\mathbb{R}$.
    A \emph{subsequence} is a sequence of the form
    \(
        (a_{n_k})_{k \in \mathbb{N}_0},
    \)
    where $(n_k)_{k \in \mathbb{N}_0}$ is a sequence of non-negative integers satisfying
    \(
        n_{k+1} > n_k \ \text{for all } k \in \mathbb{N}_0.
    \)
}
\\
\tsLem{(Limit of a sequence and its subsequence)}{
    Let $(a_n)_{n \in \mathbb{N}_0}$ be a convergent sequence with limit $L \in \mathbb{R}$.
    Then every subsequence $(a_{n_k})_{k \in \mathbb{N}_0}$ also converges to the same limit $L$.
}
\\
\tsDef{(Accumulation Point)}{
    Let $(a_n)_{n\in\mathbb{N}_0}$ be a sequence in $\mathbb{R}$.
    A number $A\in\mathbb{R}$ is called an accumulation point of the sequence if
    \[
        \forall \varepsilon>0 \;\forall N\in\mathbb{N}_0 \;\exists n\ge N \text{ such that } |a_n-A|<\varepsilon.
    \]
    Intuitively, every interval $(A-\varepsilon,A+\varepsilon)$ contains terms $a_n$ for infinitely many indices $n$.
}
\\
\tsThe{(Characterization of Accumulation Points)}{
    Let $(a_n)_{n\in\mathbb{N}_0}$ be a sequence in $\mathbb{R}$.
    A number $A\in\mathbb{R}$ is an accumulation point of the sequence if and only if there exists a convergent subsequence $(a_{n_k})_{k\in\mathbb{N}_0}$ such that
    \(
        \lim_{k\to\infty} a_{n_k}=A.
    \)
}
\\
\tsCor{(Accumulation Point of a Convergent Sequence)}{
    Every convergent sequence $(a_n)_{n\in\mathbb{N}_0}$ in $\mathbb{R}$ has exactly one accumulation point, which coincides with its limit.
}
\\
\tsCor{(Infinite Occurrence Near an Accumulation Point)}{
    Let $A$ be an accumulation point of the sequence $(a_n)_{n\in\mathbb{N}_0}$.
    Then for every $\varepsilon>0$ there exist infinitely many terms of the sequence lying in the interval
    \(
    (A-\varepsilon,A+\varepsilon).
    \)
}
\\
\tsThe{(Algebra of Limits)}{
    Assume
    \[
        \lim_{n\to\infty} a_n = K, \qquad
        \lim_{n\to\infty} b_n = L
    \]
    with $K,L\neq\pm\infty$ and let $C$ be a constant. Then
    \begin{itemize}[noitemsep]
        \item $\lim_{n\to\infty}(a_n+b_n) = K+L$
        \item $\lim_{n\to\infty}(a_n-b_n) = K-L$
        \item $\lim_{n\to\infty}(a_n b_n) = KL $
        \item $\lim_{n\to\infty}(C a_n)   = CK$
    \end{itemize}
    If $L\neq0$ and $b_n\neq0$ then
    \(
        \lim_{n\to\infty}\frac{a_n}{b_n} = \frac{K}{L}.
    \)\\
    If $K<L$ then there exists $N$ such that $a_n<b_n$ for all $n\ge N$.\\
    If there exists $N$ such that $a_n\le b_n$ for all $n\ge N$, then $K\le L$.
}
\\
\tsThe{(Important Limits)}{
    More on the last page*
    \begin{itemize}[noitemsep]
        \item
        \(
            \lim_{n\to\infty}\frac{\log n}{n}
            = 0
            = \lim_{n\to\infty}\frac{\ln n}{n}
        \)
        \item
        \(
            \lim_{n\to\infty} n^{1/n} = 1
        \)
        \item
        \(
            \lim_{n\to\infty} x^{1/n} = 1
            \qquad (x>0)
        \)
        \item
        \(
            \forall x\in\mathbb{R}:\quad
            \lim_{n\to\infty}
            \left(1+\frac{x}{n}\right)^n
            = e^x
        \)
        \item
        \(
            \forall x\in\mathbb{R}:\quad
            \lim_{n\to\infty}\frac{x^n}{n!}
            = 0
        \)
    \end{itemize}
}
\tsThe{(Sandwich Theorem)}{
    Let $(a_n)_{n\in\mathbb{N}_0}$ and $(c_n)_{n\in\mathbb{N}_0}$ be convergent sequences with
    \[
        \lim_{n\to\infty}a_n=L, \qquad
        \lim_{n\to\infty}c_n=L
    \]
    and suppose there exists $N_0$ such that
    \[
        a_n \le b_n \le c_n \quad \text{for all } n\ge N_0.
    \]
    Then $(b_n)_{n\in\mathbb{N}_0}$ is also convergent and
    \(
        \lim_{n\to\infty}b_n=L.
    \)
}
\\
\tsDef{(Bounded Sequence)}{
    A sequence $(a_n)_{n\in\mathbb{N}_0}$ is called bounded below or bounded above if the set
    \(
        M=\{a_n \mid n\in\mathbb{N}_0\}
    \)
    has a lower bound or an upper bound respectively.
    A sequence that is bounded both above and below is called a bounded sequence.
}
\\
\tsDef{(Monotone Sequence)}{
    A sequence $(a_n)_{n\in\mathbb{N}_0}$ is called:
    \begin{itemize}[noitemsep]
        \item \emph{monotone decreasing} if
        \(
            m>n \;\Rightarrow\; a_m \le a_n.
        \)

        \item \emph{monotone increasing} if
        \(
            m>n \;\Rightarrow\; a_m \ge a_n.
        \)

        \item \emph{strictly monotone decreasing} if
        \(
            m>n \;\Rightarrow\; a_m < a_n.
        \)

        \item \emph{strictly monotone increasing} if
        \(
            m>n \;\Rightarrow\; a_m > a_n.
        \)
    \end{itemize}
}
\tsThe{(Monotone Convergence Theorem)}{
    Every bounded and monotone sequence converges.
}
\\
\tsLem{(Convergent sequence = bounded)}{
    Every convergent sequence is bounded.
}
\\
\tsThe{(Monotone Convergence)}{
    Let $(a_n)_{n\in\mathbb{N}_0}$ be a sequence of real numbers.
    \begin{itemize}[noitemsep]
        \item A monotone sequence converges if and only if it is bounded.
        \item $(a_n)$ is monoton. increasing \(\Rightarrow\)
              \(
                  \lim_{n\to\infty} a_n = \sup\{a_n \mid n\in\mathbb{N}_0\}.
              \)
        \item $(a_n)$ is monoton. decreasing \(\Rightarrow\)
              \(
                  \lim_{n\to\infty} a_n = \inf\{a_n \mid n\in\mathbb{N}_0\}.
              \)
    \end{itemize}
}
\tsDef{(Limit Superior and Limit Inferior)}{
    Let $(a_n)_{n\in\mathbb{N}_0}$ be a sequence.
    \begin{itemize}[noitemsep]
        \item \emph{limit superior}:
              \(
                  \limsup_{n\to\infty} a_n
                  = \lim_{n\to\infty} \sup\{a_k \mid k \ge n\}
              \)
        \item \emph{limit inferior}:
              \(
                  \liminf_{n\to\infty} a_n
                  = \lim_{n\to\infty} \inf\{a_k \mid k \ge n\}
              \)
    \end{itemize}
}
\tsLem{(Properties of limsup and liminf)}{
    Let $(a_n)_{n\in\mathbb{N}_0}$ be a sequence. Then:
    \begin{itemize}[noitemsep]
        \item[(i)]
              \(
              \limsup_{n\to\infty} a_n
              = \inf_{n\in\mathbb{N}} \left( \sup\{a_k \mid k \ge n\} \right)
              \)

        \item[(ii)]
              \(
              \liminf_{n\to\infty} a_n
              = \sup_{n\in\mathbb{N}} \left( \inf\{a_k \mid k \ge n\} \right)
              \)

        \item[(iii)] If $(a_n)$ is convergent, then
              \[
                  \lim_{n\to\infty} a_n
                  =
                  \limsup_{n\to\infty} a_n
                  =
                  \liminf_{n\to\infty} a_n.
              \]

        \item[(iv)] A bounded sequence $(a_n)$ converges if and only if
              \[
                  \lim_{n\to\infty} a_n
                  =
                  \limsup_{n\to\infty} a_n
                  =
                  \liminf_{n\to\infty} a_n.
              \]
    \end{itemize}
}
\tsThe{(Characterization via limsup)}{
    Let $(a_n)_{n\in\mathbb{N}_0}$ be a bounded sequence and define
    \(
        A = \limsup_{n\to\infty} a_n.
    \)
    Then:
    \begin{itemize}[noitemsep]
        \item[(i)] $A$ is an accumulation point of the sequence.
        \item[(ii)] For every $\varepsilon > 0$, there are only finitely many indices $n$ such that
              \(
                  a_n > A + \varepsilon.
              \)
        \item[(iii)] For every $\varepsilon > 0$, there are infinitely many indices $n$ such that
              \(
                  A - \varepsilon < a_n < A + \varepsilon.
              \)
    \end{itemize}

    An analogous statement holds for the limit inferior.
}
\\
\tsCor{(Subsequence of a Bounded Sequence)}{
    Every bounded sequence of real numbers has at least one accumulation point and admits a convergent subsequence.
}
\\
\tsDef{(Cauchy Sequence)}{
    A sequence $(a_n)_{n\in\mathbb{N}_0} \subset \mathbb{R}$ is called a \emph{Cauchy sequence} if

    \[
        \forall \varepsilon > 0 \; \exists N \in \mathbb{N} \;\text{such that}\;
        \forall m,n \ge N:\ |a_n - a_m| < \varepsilon.
    \]
}
\tsLem{(Cauchy Sequences are Bounded)}{
    Every Cauchy sequence is bounded.
}
\\
\tsThe{(Cauchy Criterion)}{
    A sequence of real numbers converges if and only if it is a Cauchy sequence.
}